\lecture{6}{Thu 06 Feb 2025 11:02}{Product and Quotient Topologies}
A classical example is the unit circle $S^1\subseteq \mathbb{R} ^2$, where $\mathbb{R} ^2$ is given the standard topology. The unit circle can then be given the subspace topology, generated by open arcs along the circle, since every intersection of an open ball intersecting the circle will yield an open arc. We call this the standard topology on $S^1$. Some people may also use the notation $S_R^1$ for the 1-sphere with radius $R$.

We are still looking for ways to create new topologies out of old topologies. For example, the product topology.
\begin{definition}[Product Topology]\label{dfn:2.8}
	Let $X,Y$ be topological spaces with bases $\mathcal{B}_X,\mathcal{B}_Y$. Define
	\[
		\mathcal{B}_{X\times Y} \coloneqq \left\{ U\times V\mid U\in \mathcal{B}_X, V\in \mathcal{B}_Y \right\} 
	.\]
	We call the topology generated by this set the \emph{product topology} on $X\times Y$.
\end{definition}
The product space comes with natural projections $p_1 : X\times Y \to X$, $p_2 : X\times Y \to Y$ given by $(x,y)\mapsto x$ and $(x,y)\mapsto y$, respectively.

For example, consider $S_1\times S_1$. We have two projections
\[\begin{tikzcd}[row sep = 2.5em, column sep = 2.5em]
	S_1\times S_1\arrow[d]\arrow[r]&S_1\\
	S_1
\end{tikzcd}\]
Imagine placing a circle $S^1$ at each point, going along another circle. You get a taurus.
\begin{figure}[ht]
    \incfig{Taurus}
    \centering
    \caption{$S^1\times S^1$}\label{fig:Taurus}
    \centering
\end{figure}

Similarly, consider $S^1\times \mathbb{R}^+ $. This can be identified with a cylinder (stringing circles along the real line), or it can also be identified as a set with $\mathbb{R} ^2$, by using polar coordinates.

Now we will do a more indepth example. Consider $\mathbb{R} \times \mathbb{R} $ with the standard topology. Each basis element is thus the product of open intervals, so a rectangle. We showed in the first homework the collection of open rectangles generate the standard basis on $\mathbb{R} ^2$.

Now for a more formal tratment of the 2-torus. The standard topology on the Torus looks like it's generated from open rectangles from it's surface. The claim is that $S^1\times S^1$ gives this standard topology (makes sense), so we need to show this is really a basis. We can actually just prove this in a more general sense.
\begin{theorem}\label{thm:2.1}
	The basis $\mathcal{B}_{X\times Y} $ is a basis, and is independent of choice of $\mathcal{B}_X, \mathcal{B}_y$.
\end{theorem}
\begin{proof}
	For all $(x,y)\in X\times Y$, there exists $U\times V\in \mathcal{B}_{X\times Y} $ containing $(x,y)$ since $\mathcal{B}_X,\mathcal{B}_Y$ are bases, Additionally, let $U_1\times V_1,U_2\times V_2\in \mathcal{B}_{X\times Y} $, and let $(x,y)$ be contained within
	\[
		(U_1\times V_1)\cap (U_2\times V_2)
	.\]
	Since $\mathcal{B}_X,\mathcal{B}_Y$ are bases, there exist $U_3\subseteq U_1 \cap U_2$ and $V_3 \subseteq V_1 \cap V_2$ satisfying this stuff, so we have that
	\[
		(x,y)\in U_3 \times V_3 \subseteq (U_1\times V_1)\cap (U_2\times V_2)
	.\]
	Therefore, this is a basis. Finally, we must show the choice of basis is irrelevant. Let $\mathcal{B}_X,\mathcal{B}_X'$ be bases $X$, and let $\mathcal{B}_Y,\mathcal{B}_Y'$ be bases for $Y$. It suffices to show that
	\[
		\mathcal{T}_{\mathcal{B}_{X\times Y} } \subseteq \mathcal{T}_{\mathcal{B}_{X\times Y}'} 
	.\]
	This is because we can use the same statement in the other direction to show equality.

	Let
	\[
		\bigcup_{i\in I} U_i\times V_i \in \mathcal{B}_{X\times Y} 
	,\]
	where $U_i\in \mathcal{B}_X$, $V_i\in \mathcal{B}_Y$ for all $i\in I$. Fix $x\in U_i$. By properties of basis, there exists $U_{i,x}'\in \mathcal{B}_X'$ such that
	\[
		x\in U_{i,x}' \subseteq U_i
	.\]
	We can thus write
	\[
		U_i = \bigcup_{x\in U_i} U_{i,x}'
	.\]
	We can apply the same logic to the $V_i$s, showing that it can be written as a union of elements in $\mathcal{B}_Y'$. Therefore, we are done.
\end{proof}

We can create yet another topology, called the quotient topology.
\begin{definition}[Quotient Topology]\label{dfn:2.9}
	Let $A$ be a nonempty set and let $X$ be a topological space. Let $p : X \twoheadrightarrow A$ be surjective. Define the quotient topology on $A$ such that $U\subseteq A$ if open if and only if $p ^{-1} (U)$ is open in $X$.
\end{definition}
\begin{proposition}\label{prp:2.9}
	Let $A$ be a nonempty set and $X$ a topological space, and let $p : X \twoheadrightarrow A$ a surjection. The quotient topology on $A$ by $p$ defines a topology on $A$.
\end{proposition}
\begin{proof}
	Let $\mathcal{T}_A$ be the collection of all subsets of $A$ with open preimages. We first want to show $\varnothing ,A\in \mathcal{T}_A$. Since $p$ is surjective, $p ^{-1} (\varnothing )=\varnothing $. Additionally, since $p$ is surjective, $p ^{-1} (A)=X$. Since $\varnothing ,X$ are open in $X$, we have $\varnothing ,A\in \mathcal{T}_A$ are open.

	Now let $U_1, \ldots , U_k$ be open in $A$. Note that since $p ^{-1} (U_i)$ is open in $X$, and finite intersection of open sets is open,
	\[
		\bigcap_{i=1}^k p ^{-1} (U_i)
	\]
	is open in $X$. Therefore,
	\[
		\bigcap_{i=1}^k U_i \in \mathcal{T}_A
	.\]
	Finally, let $\left\{U_i \right\}_{i\in I} \subseteq \mathcal{T}_A$. Then by the same logic,
	\[
		\bigcup_{i\in I} p ^{-1}(U_i)
	\]
	is open in $X$. It follows that
	\[
		p\left( \bigcup_{i\in I} p ^{-1} (U_i) \right) =\bigcup_{i\in I}U_i \in \mathcal{T}_A
	.\]
\end{proof}

For example, consider $\mathbb{R} $ and $S^1$, with the map $\mathbb{R} \to S^1$ given by $t\mapsto \exp (2\pi it)$. The prof is lost trying to explain what this topology is. We can identify $[0,1)$ with $S^1$ by the map $p$. Ok so an open set is basically the image of an interval, which is going to be an open arc, so the quotient topology is the standard topology on $S^1$.

Given any relation $\sim $ on $X$, consider the quotient map in $\mathbf{Set} $ given by $x\mapsto [x]$. Given that $(X,\mathcal{T} )$ is a topological space, we can then give $X /\tildesim $ the quotient topology by the quotient map. An application of this is taking the relation $(0,y)\sim (1,y)$ on $[0,1]\times [0,1]$. This glues together two sides of the rectangle, creating the cylinder $S^1 \times [0,1]$.
